% Volume VII: Deep Learning Mathematics
% Continuity note for Chapter 42.
% Previous dependencies: Chapters 5-7 introduced geometry, projections, PCA, tensors, and matrix calculus; Chapters 18 and 31-38 introduced information, learning, representations, neural networks, and backpropagation; Chapters 39-41 introduced training dynamics, architectures, and generative models.
% New durable concepts: representation matrices, distances, angles, covariance, linear probes, RSA/CKA intuition, adversarial perturbations, robust risk, dual-norm attacks, margins, saliency, integrated gradients, probes, activation patching, causal caveats, scaling laws, and emergence measurement caveats.
% Forward hooks: Chapters 47-50 reuse representation geometry for neural coding, population manifolds, attractors, and predictive coding; Chapter 52 compares biological and artificial intelligence with these caveats in place.
\section{Chapter 42: Representation Geometry, Robustness, and Interpretability}
\label{ch:representation-geometry-robustness-interpretability}

Deep networks do not only output predictions. They transform inputs into internal representations. Those representations have geometry: distances, directions, clusters, manifolds, and low-dimensional projections. They also have vulnerabilities: small mathematical perturbations can change predictions, and the features that support a decision may not match human concepts. Interpretability tries to connect internal computation to understandable mechanisms, but it must distinguish correlation, decodability, and causal use.

\paragraph{Chapter problem.}
How can we analyze learned representations geometrically, test their robustness, interpret their mechanisms, and reason about scaling without overclaiming what the model or biology proves?

\paragraph{Section map.}
Section 42.1 studies representation geometry through representation matrices, distances, covariance, linear probes, and comparisons to neural population geometry. Section 42.2 defines robustness and adversarial examples using perturbation sets, margins, and robust risk. Section 42.3 studies interpretability methods such as saliency, probes, and causal interventions, with their failure modes. Section 42.4 introduces scaling laws and emergence caveats.

\subsection{42.1 Geometry of learned representations}

\paragraph{Guiding question.}
Once a network maps examples to hidden vectors, what geometric questions can we ask about those vectors?

\paragraph{Beginner-facing intuition.}
A representation is a coordinate system learned for a task. At layer \(\ell\), a network maps each input \(x_i\) to a hidden vector
\[
h_i=h_\ell(x_i)\in\mathbb{R}^d.
\]
If similar inputs have nearby hidden vectors, if classes become linearly separable, or if important variables align with low-dimensional directions, then the representation has useful geometry.

Geometry is not the same as interpretation. Distances, angles, and clusters are mathematical properties of a representation matrix. Their scientific meaning depends on the data, task, and validation method.

\paragraph{Formal definitions and notation.}
For \(n\) examples, collect hidden vectors as rows of a representation matrix
\[
H=
\begin{bmatrix}
h_1^\top\\
\vdots\\
h_n^\top
\end{bmatrix}
\in\mathbb{R}^{n\times d}.
\]
Pairwise squared distances are
\[
D_{ij}=\|h_i-h_j\|_2^2.
\]
The centered representation is
\[
\widetilde H=H-\mathbf{1}\bar h^\top,\qquad
\bar h=\frac1n\sum_{i=1}^{n}h_i,
\]
and its empirical covariance is
\[
C=\frac1n\widetilde H^\top\widetilde H.
\]
A linear probe predicts a target \(y_i\) from \(h_i\) using a simple model such as
\[
\widehat y_i=w^\top h_i+b.
\]
Representational similarity analysis (RSA) compares pairwise distance or similarity matrices across models or brains. Centered kernel alignment (CKA) compares representation Gram matrices after centering and normalization.

\paragraph{Core mechanism: PCA directions of a representation.}
Principal component analysis of the centered representation diagonalizes the covariance
\[
C v_k=\lambda_k v_k,
\qquad
\lambda_1\geq\lambda_2\geq\cdots\geq0.
\]
The centered score of example \(i\) along direction \(v_k\), equivalently the \(i\)-th row of \(\widetilde H\) multiplied by \(v_k\), is
\[
u_{ik}=(h_i-\bar h)^\top v_k.
\]
Large eigenvalues indicate directions along which hidden vectors vary strongly. If most variance lies in a few directions, the representation is approximately low-dimensional in the variance sense. This is a geometric statement; it does not by itself say that the directions are causal or semantically meaningful.

\paragraph{Concrete example.}
Suppose four examples have two-dimensional hidden vectors
\[
h_1=(0,0),\quad h_2=(0,1),\quad h_3=(3,0),\quad h_4=(3,1).
\]
The first coordinate separates the pairs \(\{h_1,h_2\}\) and \(\{h_3,h_4\}\). A linear classifier with \(w=(1,0)\) and threshold \(1.5\) separates left from right. The second coordinate varies within each pair. A probe can detect which labels are linearly accessible from this representation.

\paragraph{Computational neuroscience example.}
Neural population geometry uses similar objects. If \(H\) contains trial-by-neuron firing-rate vectors, distances and low-dimensional projections describe how neural states vary across stimuli or behaviors. Comparing a model representation to a neural representation via RSA can reveal similar geometry, but it does not prove that the model and brain use the same algorithm.

\paragraph{AI/ML example.}
In a classifier, early layers may preserve local input geometry while later layers spread classes apart. In a language model, hidden states may make syntactic number, topic, or token identity linearly decodable. A linear probe tests accessibility of information, but high probe accuracy does not prove the main model uses that information for its output.

\paragraph{Common misconceptions and failure modes.}
Low-dimensional visualization can create misleading clusters because projection discards information. A linear probe measures decodability, not necessarily causal use. Similar RSA matrices can arise from different mechanisms. Distance in hidden space depends on scaling, normalization, and layer choice. A representation that matches neural data on one benchmark may fail on another stimulus set or task.

\paragraph{Exercises.}
\begin{enumerate}[label=\textbf{42.1.\arabic*.}, leftmargin=*]
\item \textbf{Mechanical.} For \(h_1=(1,0)\), \(h_2=(0,1)\), and \(h_3=(1,1)\), compute all pairwise squared distances.
\item \textbf{Proof/derivation.} Show that the empirical covariance \(C=(1/n)\widetilde H^\top\widetilde H\) is positive semidefinite.
\item \textbf{Numerical/computational.} Given hidden vectors \((0,0),(2,0),(0,1),(2,1)\), find a linear rule that separates examples with first coordinate \(0\) from examples with first coordinate \(2\).
\item \textbf{Interpretation.} What does high linear-probe accuracy show, and what does it not show, about a feature in a model representation?
\item \textbf{Misconception/debugging.} A student sees two clusters in a two-dimensional PCA plot and concludes that the network has exactly two internal states. Explain why this conclusion is too strong.
\end{enumerate}

\subsection{42.2 Robustness and adversarial examples}

\paragraph{Guiding question.}
How can a model be accurate on ordinary test data but fragile to small, carefully chosen perturbations?

\paragraph{Beginner-facing intuition.}
Robustness asks whether predictions remain stable under allowed changes. The allowed changes must be defined mathematically: an \(\ell_\infty\) pixel perturbation, an \(\ell_2\) sensor noise ball, a rotation, a translation, a distribution shift, or a biologically plausible stimulus change. An adversarial example is an input modified within an allowed set to make the model fail.

Small in a norm does not always mean imperceptible, and imperceptible to a human does not always mean small in a simple norm. Robustness is therefore both a mathematical and modeling question.

\paragraph{Formal definitions and notation.}
Let \(f_\theta:\mathcal X\to\mathcal Y\) be a classifier, let \(\ell(f_\theta(x),y)\) be a loss, and let
\[
\Delta_\epsilon=\{\delta:\|\delta\|\leq\epsilon\}
\]
be an allowed perturbation set. The robust loss at \((x,y)\) is
\[
\boxed{
L_{\mathrm{rob}}(x,y)
=\max_{\delta\in\Delta_\epsilon}
\ell(f_\theta(x+\delta),y).
}
\]
The robust risk is the expectation of this quantity over the data distribution. For a binary linear classifier
\[
f(x)=\operatorname{sign}(w^\top x+b),\qquad y\in\{-1,1\},
\]
the signed margin is
\[
m=y(w^\top x+b).
\]
The example is correctly classified when \(m>0\).

\paragraph{Core derivation: FGSM from linearizing the loss.}
For a differentiable loss, linearize around \(x\):
\[
\ell(f_\theta(x+\delta),y)
\approx
\ell(f_\theta(x),y)+\nabla_x\ell(f_\theta(x),y)^\top\delta.
\]
Under an \(\ell_\infty\) constraint \(\|\delta\|_\infty\leq\epsilon\), the linear term is maximized by choosing each coordinate to match the sign of the gradient:
\[
\boxed{
\delta_{\mathrm{FGSM}}
=\epsilon\,\operatorname{sign}(\nabla_x\ell(f_\theta(x),y)).
}
\]
This is the fast gradient sign method. More iterative attacks, such as projected gradient descent, repeatedly take gradient steps and project back into the allowed perturbation set.

For a linear classifier and a norm \(\|\cdot\|\), Holder's inequality gives the worst margin decrease
\[
\max_{\|\delta\|\leq\epsilon}[-y w^\top\delta]
=\epsilon\|w\|_*,
\]
where \(\|\cdot\|_*\) is the dual norm. Thus the classifier is certifiably robust for that example if
\[
y(w^\top x+b)>\epsilon\|w\|_*.
\]

\paragraph{Concrete example.}
Let \(w=(2,1)\), \(b=0\), \(x=(0.4,0.3)\), and \(y=1\). The margin is
\[
m=2(0.4)+1(0.3)=1.1.
\]
For \(\ell_\infty\) perturbations of size \(\epsilon=0.2\), the dual norm is \(\|w\|_1=3\), so the worst margin decrease is \(0.6\). Since \(1.1>0.6\), this linear classifier is robust for this example under that perturbation size. If \(\epsilon=0.5\), the worst decrease is \(1.5\), so this certificate fails.

\paragraph{Computational neuroscience example.}
Sensory systems face noise, occlusion, eye movements, and changes in lighting. Robustness analysis can ask whether a neural encoding or decoding model is stable under plausible perturbations of stimuli or recordings. But adversarial pixel perturbations designed for an artificial classifier are not automatically plausible biological perturbations. The perturbation set must match the scientific question.

\paragraph{AI/ML example.}
Image classifiers can be fooled by small \(\ell_\infty\) perturbations that preserve human-recognizable content. Speech and language models have analogous vulnerabilities under audio perturbations or token substitutions. Adversarial training replaces ordinary loss with an approximation to robust loss, improving robustness to the chosen attack class but often increasing compute and sometimes reducing clean accuracy.

\paragraph{Common misconceptions and failure modes.}
Robust to one perturbation set does not mean robust to all shifts. A failed attack is not a proof of robustness unless the attack is complete for the threat model. Human imperceptibility is not identical to \(\ell_p\)-smallness. Adversarial examples do not prove a model has learned nothing; they reveal a mismatch between learned decision boundaries and the chosen robustness criterion.

\paragraph{Exercises.}
\begin{enumerate}[label=\textbf{42.2.\arabic*.}, leftmargin=*]
\item \textbf{Mechanical.} For gradient \(g=(-2,0.5,3)\) and \(\epsilon=0.1\), compute the FGSM perturbation under an \(\ell_\infty\) constraint.
\item \textbf{Proof/derivation.} Use coordinate-wise maximization to prove that \(\max_{\|\delta\|_\infty\leq\epsilon}g^\top\delta=\epsilon\|g\|_1\).
\item \textbf{Numerical/computational.} For \(w=(1,-2)\), \(b=0.1\), \(x=(2,1)\), and \(y=1\), compute the margin and determine whether the example is certified robust to \(\ell_\infty\) perturbations with \(\epsilon=0.2\).
\item \textbf{Interpretation.} Why must a neuroscience robustness claim specify whether perturbations are pixel noise, stimulus transformations, neural recording noise, or behavioral context changes?
\item \textbf{Misconception/debugging.} A student says, ``The model resisted FGSM, so it is adversarially robust.'' Explain why stronger attacks, threat-model definition, and certificates matter.
\end{enumerate}

\subsection{42.3 Interpretability}

\paragraph{Guiding question.}
What kinds of evidence can connect a model's internal variables to features, decisions, or mechanisms?

\paragraph{Beginner-facing intuition.}
Interpretability is not one method. A saliency map asks which input coordinates locally affect a score. A probe asks what information is decodable from a representation. A circuit analysis asks which components transform information into a behavior. A causal intervention asks what changes when an activation, feature, or connection is modified.

The key distinction is between \emph{correlation}, \emph{decodability}, and \emph{causal use}. A feature may be present in a representation but unused by the output. A gradient may be large locally but unstable under small changes. A component may matter in one context but not another.

\paragraph{Formal definitions and notation.}
Let \(s_c(x)\) be the scalar score or logit for class \(c\). A gradient saliency map is
\[
S(x)=\left|\nabla_x s_c(x)\right|,
\]
with one saliency value per input coordinate. Integrated gradients compare input \(x\) to a baseline \(x'\):
\[
\operatorname{IG}_i(x)
=(x_i-x_i')\int_0^1
\frac{\partial s_c(x'+\alpha(x-x'))}{\partial x_i}\,d\alpha.
\]
A linear probe for representation \(h(x)\) fits
\[
\widehat y=w^\top h(x)+b
\]
to predict a target attribute. An activation intervention replaces a hidden activation \(h_j\) with another value \(\widetilde h_j\) and measures the change in output or loss.

\paragraph{Core mechanism: saliency for a linear model.}
For a linear score
\[
s_c(x)=w^\top x+b,
\]
the gradient is
\[
\nabla_x s_c(x)=w.
\]
Thus saliency is constant across inputs and equals the absolute weight magnitudes. If \(w_i\) is large, changing coordinate \(x_i\) locally changes the score strongly. For nonlinear networks, the gradient depends on \(x\), so saliency is a local linear explanation:
\[
s_c(x+\delta)\approx s_c(x)+\nabla_x s_c(x)^\top\delta.
\]
This local nature is both the strength and weakness of gradient explanations.

\paragraph{Concrete example.}
Let
\[
s(x)=3x_1-x_2+0.5x_3.
\]
Then
\[
\nabla_x s=(3,-1,0.5),
\]
and the saliency magnitudes are \((3,1,0.5)\). Coordinate \(x_1\) has the largest local effect on the score. This does not say \(x_1\) is semantically meaningful; it says the score is locally sensitive to that coordinate in this coordinate system.

\paragraph{Computational neuroscience example.}
Neuroscience has long used perturbations, lesions, stimulation, and decoding analyses. Model interpretability borrows analogous logic: decoding a variable from neural activity resembles probing; inactivating a region resembles intervention. But the analogy has limits. A probe trained on a representation does not prove the system uses that information, just as decoding a stimulus from a brain area does not by itself prove that area drives behavior.

\paragraph{AI/ML example.}
In a vision model, saliency may highlight pixels that affect a class score. In a language model, activation patching can replace an activation from one prompt with the corresponding activation from another prompt to test whether a component mediates a behavior. In both cases, rigorous interpretation requires controls: baselines, randomization tests, alternative explanations, and out-of-distribution checks.

\paragraph{Common misconceptions and failure modes.}
A pretty heatmap is not an explanation unless the method and question are defined. Gradients can saturate or be noisy. Probes can be too powerful and extract information the main model does not use. Ablating a component can create off-distribution activations, so the measured effect may combine mechanism and damage. Interpretability claims should state the level of evidence: sensitivity, decodability, necessity, sufficiency, or causal mediation.

\paragraph{Exercises.}
\begin{enumerate}[label=\textbf{42.3.\arabic*.}, leftmargin=*]
\item \textbf{Mechanical.} For \(s(x)=2x_1+4x_2-x_3\), compute the gradient saliency magnitudes.
\item \textbf{Proof/derivation.} Show that for \(s(x)=w^\top x+b\), the integrated gradients from baseline \(0\) satisfy \(\operatorname{IG}_i(x)=x_i w_i\).
\item \textbf{Numerical/computational.} A probe predicts a binary attribute from hidden vectors with \(95\%\) accuracy, but a control probe on shuffled labels gets \(90\%\). What concern does this raise?
\item \textbf{Interpretation.} Explain the difference between ``a feature is linearly decodable from \(h\)'' and ``the model causally uses that feature.''
\item \textbf{Misconception/debugging.} A saliency map highlights the background of an image. Give two possible explanations other than ``the model understands the background concept.''
\end{enumerate}

\subsection{42.4 Scaling and emergence}

\paragraph{Guiding question.}
How do model size, data, and compute change performance, and what should we mean by an emergent capability?

\paragraph{Beginner-facing intuition.}
Scaling studies ask how loss or performance changes as resources grow. The resources may be parameter count \(N\), dataset size \(D\), compute \(C\), context length, or environment interactions. Empirically, many losses follow approximate power laws over some range. This makes scaling predictable enough to guide engineering, but not magical: the fitted law depends on data quality, architecture, optimization, and metric.

Emergence refers to capabilities that appear sharply as scale increases. Sometimes this reflects a real qualitative change in behavior. Sometimes it reflects a thresholded metric: a smoothly improving probability becomes a sudden jump in measured accuracy once it crosses a cutoff.

\paragraph{Formal definitions and notation.}
A simple one-variable scaling law for loss as a function of parameter count \(N\) is
\[
\boxed{
L(N)=L_\infty+aN^{-\alpha},
}
\]
where \(L_\infty\) is an irreducible loss floor, \(a>0\), and \(\alpha>0\). A two-resource form is
\[
L(N,D)=L_\infty+aN^{-\alpha}+bD^{-\beta},
\]
where \(D\) is data size. On a log-log plot, if \(L(N)-L_\infty=aN^{-\alpha}\), then
\[
\log(L(N)-L_\infty)=\log a-\alpha\log N,
\]
so the slope is \(-\alpha\).

Let \(M(N)\) be a measured capability score. A thresholded metric may be
\[
A(N)=\mathbf{1}\{M(N)\geq \tau\}.
\]
Even if \(M(N)\) changes smoothly with \(N\), \(A(N)\) jumps at the threshold.

\paragraph{Core mechanism: fitting a power-law slope.}
Suppose \(L_\infty=1\), and measured losses are
\[
L(100)=3,\qquad L(400)=2.
\]
Then the excess losses are
\[
L(100)-1=2,\qquad L(400)-1=1.
\]
Assuming \(L(N)-1=aN^{-\alpha}\),
\[
\frac{1}{2}=\frac{a400^{-\alpha}}{a100^{-\alpha}}
=\left(\frac{400}{100}\right)^{-\alpha}
=4^{-\alpha}.
\]
Thus \(\alpha=1/2\). This calculation is useful only if the loss floor and power-law range are credible.

\paragraph{Concrete example.}
Let
\[
L(N)=1+10N^{-1/2}.
\]
At \(N=100\), \(L=2\). At \(N=400\), \(L=1.5\). At \(N=1600\), \(L=1.25\). Increasing scale helps, but with diminishing returns. To halve the excess loss again, \(N\) must grow by another factor of \(4\).

\paragraph{Computational neuroscience example.}
Scaling ideas can be used when fitting models to neural data: more neurons, more trials, larger models, or longer recordings can improve predictive performance. But biological comparison requires caution. A larger artificial model matching more neural variance does not imply that the brain implements the same architecture, and neural data may have measurement noise floors that look like \(L_\infty\).

\paragraph{AI/ML example.}
Large language models often show smoother improvements in next-token loss than in downstream benchmark accuracy. A benchmark may appear emergent when a continuous improvement in probability finally crosses the threshold required to answer items correctly. Other behaviors may involve genuine new combinations of skills enabled by scale, data, and training. The mathematical lesson is to inspect the underlying metric before declaring a phase transition.

\paragraph{Common misconceptions and failure modes.}
Scaling laws are empirical fits over a range, not universal guarantees. Parameter count alone is not scale; data, compute, architecture, optimization, and evaluation all matter. Emergence should not be inferred from one thresholded plot without checking smoother metrics. Bigger models can amplify data biases and vulnerabilities. Biological analogies based on scale should specify what is being scaled in the model and what would correspond to it in neural systems.

\paragraph{Exercises.}
\begin{enumerate}[label=\textbf{42.4.\arabic*.}, leftmargin=*]
\item \textbf{Mechanical.} For \(L(N)=0.5+4N^{-1/2}\), compute \(L(100)\) and \(L(400)\).
\item \textbf{Proof/derivation.} Show that \(L(N)-L_\infty=aN^{-\alpha}\) becomes a straight line with slope \(-\alpha\) on a log-log plot.
\item \textbf{Numerical/computational.} If excess loss drops from \(8\) to \(2\) when \(N\) increases from \(10^3\) to \(10^5\), estimate \(\alpha\).
\item \textbf{Interpretation.} Explain how a thresholded benchmark can make a smooth capability curve look emergent.
\item \textbf{Misconception/debugging.} A student says, ``The larger model matches neural data better, so it is more brain-like.'' Give two alternative explanations and one additional test you would want.
\end{enumerate}

\paragraph{Closing bridge.}
This chapter completes the deep-learning mathematics volume by separating useful analysis from overinterpretation. Representation geometry gives measurable structure; robustness defines stability under specified perturbations; interpretability methods provide evidence at different strengths; scaling laws summarize empirical trends with caveats. The next volume returns to biological neural systems, where these tools must be connected to biophysics, spikes, synapses, neural codes, and circuits with explicit mathematical assumptions.

\clearpage
